\documentclass[11pt]{article}
\usepackage[margin=1in]{geometry}
\usepackage{booktabs}
\usepackage{hyperref}
\usepackage{microtype}

\title{Does Wrong-Answer Conditioning Change Educational Evidence Retrieval?\\A Controlled SciQ Study}
\author{Devis Wawan Saputra}
\date{}

\begin{document}
\maketitle

\begin{abstract}
Misconception-aware retrieval systems often assume that an incorrect response can improve evidence selection, but appending any text can also alter lexical retrieval. We evaluate question-only BM25, observed wrong-answer conditioning, a shuffled wrong-answer negative control, and an oracle-informed gold-answer sensitivity condition on a frozen SciQ test split. Rank metrics and question-block bootstrap intervals separate the primary retrieval effect from a generic query-expansion explanation.
\end{abstract}

\section{Research question}
Does adding an observed wrong-answer option change retrieval of a science question's support passage, and is that effect distinguishable from adding unrelated distractor-like text?

\section{Data}
We use the SciQ test split \cite{welbl2017sciq} pinned to a fixed dataset revision and SHA-256. Only items with non-empty support passages enter the study. SciQ distractors are treated as wrong-answer proxies, not validated misconception labels.

\section{Retrieval conditions}
BM25 \cite{robertson2009bm25} uses $k_1=1.5$, $b=0.75$, lowercased Unicode word tokens, and retrieval depth five. Conditions are question-only; question plus observed distractor; question plus a shuffled distractor from a different eligible question; and question plus observed distractor plus gold correct answer. The shuffled condition preserves the marginal distractor-text distribution while breaking question--answer alignment. The gold-answer condition is oracle-informed sensitivity analysis, not a mathematical upper bound.

\section{Evaluation}
We report MRR, Recall@1/3/5 and nDCG@5. Since each question yields three distractor cases, uncertainty is resampled by question block. We report observed-wrong-answer versus question-only, shuffled-control versus question-only, and observed-wrong-answer versus shuffled-control contrasts.

\section{Generated empirical results}
This section is generated by \texttt{scripts/run\_sciq\_study.py}; numerical values should not be hand-edited.

Frozen SciQ revision: \texttt{2c94ad3e1aafab77146f384e23536f97a4849815}; eligible questions: 884; wrong-answer proxy cases: 2652.

\begin{table}[htbp]
\centering
\small
\begin{tabular}{lrrrrr}
\toprule
Condition & MRR & R@1 & R@3 & R@5 & nDCG@5 \\
\midrule
Question only & 0.9472 & 0.9253 & 0.9672 & 0.9740 & 0.9540 \\
Observed wrong answer & 0.9433 & 0.9197 & 0.9672 & 0.9732 & 0.9509 \\
Shuffled wrong-answer control & 0.9447 & 0.9223 & 0.9664 & 0.9713 & 0.9516 \\
Oracle-informed corrective & 0.9612 & 0.9453 & 0.9766 & 0.9815 & 0.9664 \\
\bottomrule
\end{tabular}
\caption{SciQ support-passage retrieval under the four frozen query conditions.}
\end{table}

Observed wrong answer minus question-only: mean MRR delta -0.0039, question-block bootstrap 95\% interval [-0.0077, -0.0002].
Shuffled control minus question-only: mean MRR delta -0.0024, question-block bootstrap 95\% interval [-0.0044, -0.0005].
Observed wrong answer minus shuffled control: mean MRR delta -0.0014, question-block bootstrap 95\% interval [-0.0057, 0.0027].

The shuffled control preserves the marginal distractor-text distribution while breaking the question--answer relation. It is a lexical-expansion negative control, not a simulated learner response.


\section{Validity and responsible-use boundary}
The benchmark measures lexical retrieval of SciQ support passages, not misconception diagnosis, tutoring quality or learning gain. The shuffled control helps diagnose generic query-expansion effects but does not model learner cognition. The separate tutoring prototype remains unvalidated for learner-facing use.

\bibliographystyle{plain}
\bibliography{references}
\end{document}
